% Part: model-theory
% Chapter: models-of-arithmetic
% Section: non-standard-models

\documentclass[../../../include/open-logic-section]{subfiles}

\begin{document}

\olfileid[fr]{mod}{mar}{nst}
\section{Modèles non standards}

\begin{explain}
Une !!{structure} pour $\Lang{L_A}$ est dite standard
si elle est isomorphe à~$\Struct{N}$; sinon, on
la dit non standard.
\end{explain}

\begin{defn}
Une !!{structure}~$\Struct{M}$ pour $\Lang{L_A}$
est \emph{non standard} si elle n'est pas isomorphe
à~$\Struct{N}$. Les !!{element}s $x\in\Domain{M}$
qui sont égaux à $\Value{\num{n}}{M}$ pour un certain
$n\in\Nat$ sont les \emph{nombres standards}
de~$\Struct{M}$; les autres sont ses
\emph{nombres non standards}.
\end{defn}

\begin{explain}
D'après la \olref[stm]{prop:standard-domain}, une
!!{structure} standard ne contient que des nombres
standards. La réciproque est fausse pour une
!!{structure} arbitraire, comme le montre l'exercice
suivant cette proposition : une !!{structure} non
standard peut avoir tous ses éléments nommés par
des numéraux. Si elle est en outre un modèle
de~$\Th{Q}$, \olref[stm]{prop:thq-standard} exclut
ce cas. En revanche, la présence d'un nombre
non standard suffit toujours à montrer que
la !!{structure} est non standard.\footnote{Correction éditoriale : la source affirme aussi qu'une structure non standard contient nécessairement un élément non nommé par un numéral. Cette réciproque est fausse pour les structures arbitraires ; elle devient vraie pour les modèles de $\Th{Q}$ par la proposition citée.}
\end{explain}

\begin{prop}
Si une !!{structure}~$\Struct{M}$ pour $\Lang{L_A}$
contient un nombre non standard, elle est non standard.
\end{prop}

\begin{proof}
Supposons au contraire que $\Struct{M}$ soit
standard et contienne un nombre non standard~$x$.
Soit $g\colon\Nat\to\Domain{M}$ un isomorphisme.
Par récurrence sur~$n$,
$g(\Value{\num{n}}{N})=\Value{\num{n}}{M}$ :
$g$ envoie donc les nombres standards de~$\Struct{N}$
sur les nombres standards de~$\Struct{M}$. Comme
tout élément de~$\Nat$ est la valeur d'un numéral,
$x$ n'appartient pas à l'image de~$g$. Cela
contredit la surjectivité de~$g$.
\end{proof}

\begin{prob}
Rappelons que $\Th{Q}$ contient les axiomes
\begin{align*}
& \lforall[x][\lforall[y][(\eq[x'][y'] \lif \eq[x][y])]]
  \tag{$!Q_1$}\\
& \lforall[x][\eq/[\Obj 0][x']] \tag{$!Q_2$}\\
& \lforall[x][(\eq[x][\Obj 0]\lor
  \lexists[y][\eq[x][y']])] \tag{$!Q_3$}.
\end{align*}
Donner trois !!{structure}s
$\Struct{M_1}$, $\Struct{M_2}$ et $\Struct{M_3}$
telles que
\begin{enumerate}
\item $\Sat{M_1}{!Q_1}$, $\Sat{M_1}{!Q_2}$,
  $\Sat/{M_1}{!Q_3}$;
\item $\Sat{M_2}{!Q_1}$, $\Sat/{M_2}{!Q_2}$,
  $\Sat{M_2}{!Q_3}$;
\item $\Sat/{M_3}{!Q_1}$, $\Sat{M_3}{!Q_2}$,
  $\Sat{M_3}{!Q_3}$.
\end{enumerate}
Il suffit de préciser, pour chaque~$i$, les valeurs
de $\Assign{\Obj{0}}{M_i}$ et
de~$\Assign{\prime}{M_i}$.
\end{prob}

\begin{explain}
Il est facile de décrire une !!{structure} non standard
pour $\Lang{L_A}$ : prenons $\Int$ comme domaine et
interprétons tous les symboles non logiques de la façon
habituelle. Aucun entier négatif n'est la valeur
d'un numéral~$\num{n}$; cette !!{structure} est donc
non standard. Elle n'est cependant pas un modèle
d'une théorie arithmétique qui rende vrais les
mêmes !!{sentence}s que~$\Struct{N}$.
Par exemple,
$\lforall[x][\eq/[x'][\Obj{0}]]$
y est faux, car $(-1)'=0$. La compacité permet
néanmoins de construire des modèles non standards
de l'arithmétique.
\end{explain}

\begin{prop}
Soit
$\Th{TA}=\Setabs{!A}{\Sat{N}{!A}}$
la théorie de~$\Struct{N}$.
La théorie $\Th{TA}$ possède un modèle non standard
!!{enumerable}.
\end{prop}

\begin{proof}
Ajoutons au !!{language}~$\Lang{L_A}$ une nouvelle
!!{constant}~$c$ et considérons
\[
  \Gamma=\Th{TA}\cup
  \{\eq/[c][\num{0}],\eq/[c][\num{1}],
    \eq/[c][\num{2}],\dots\}.
\]
Dans tout modèle~$\Struct{M^c}$ de~$\Gamma$,
l'élément $x=\Assign{c}{M^c}$ est distinct de
$\Value{\num{n}}{M^c}$ pour tout~$n$.
Le réduit~$\Struct{M}$ de $\Struct{M^c}$ au
!!{language}~$\Lang{L_A}$ contient toujours
cet élément non standard et satisfait~$\Th{TA}$,
car $c$ n'apparaît dans aucun de ses énoncés.\footnote{Correction éditoriale : la constante nouvelle $c$ est interprétée dans $\Struct{M^c}$ ; elle n'appartient pas au langage du réduit $\Struct{M}$, contrairement à la notation de la source.}

Montrons donc que $\Gamma$ a un modèle.
Soit $\Gamma_0\subseteq\Gamma$ fini. Si
$\Gamma_0$ contient des inégalités
$\eq/[c][\num{n}]$, choisissons~$k$ plus
grand que tous les indices~$n$ qui y figurent;
s'il n'en contient aucune, choisissons par exemple
$k=0$.\footnote{Complément éditorial : la source choisit le plus grand indice d'une inégalité ; nous traitons aussi les sous-ensembles finis qui n'en contiennent aucune.} Développons $\Struct{N}$ en
$\Struct{N_k}$ en posant
$\Assign{c}{N_k}=k+1$.
Tous les énoncés de
$\Gamma_0\cap\Th{TA}$ y sont vrais, car
$c$ n'y apparaît pas. Les inégalités de
$\Gamma_0$ y sont vraies aussi, puisque leurs
indices sont inférieurs à~$k+1$.
Ainsi, tout sous-ensemble fini de~$\Gamma$
est satisfaisable; par compacité, $\Gamma$
possède un modèle. Enfin, le langage étendu
est dénombrable et tout modèle de~$\Gamma$
est infini, puisque les numéraux y désignent
des éléments distincts. Le théorème de
L\"owenheim--Skolem descendant fournit donc
un modèle dénombrable de~$\Gamma$; son réduit
à $\Lang{L_A}$ est le modèle annoncé.\footnote{Complément éditorial : la preuve source s'arrête à l'existence d'un modèle. Le recours au théorème descendant fournit la cardinalité dénombrable annoncée ; le modèle du langage étendu conserve l'interprétation de $c$ et donc l'élément non standard.}
\end{proof}

\end{document}
